\chapter{Literature Review} \label{ch:litreview}
    This chapter reviews the strands of prior work that underpin the constellation-design and optimization choices used in the remainder of the dissertation. First, it summarizes analytic constellation families and search-based design methods, emphasizing results that inform coverage, revisit, and maintainability. Next, it reviews stochastic optimization formulations and tools that treat uncertainty as a first-class design driver. Throughout, the chapter highlights gaps that motivate the dissertation's approach to uncertainty-aware, operations-informed constellation design for debris-remediation tasking.
    \section{Constellation Optimization} \label{se:consteloptim}
    
    Constellation design is often framed around (i) analytic families with strong geometric structure that admit closed-form coverage reasoning and simple station-keeping, and (ii) search-based methods that tune architectures against multi-criteria performance measures \cite{lang1998comp}. Walker constellations remain a common analytic baseline \cite{walker1984,Walker1985,lang1998comp}. In Walker's ``\(T/P/F\)'' notation, \(T\) satellites are distributed across \(P\) orbital planes with evenly spaced \gls{raan}, while \(F\) sets inter-plane phasing \cite{lang1998comp}. Real systems anchor this taxonomy: Galileo's reference constellation uses a \(56^{\circ}\) inclination with three primary plane groups in \gls{raan} \cite{GalileoSISICD2025}, and FCC filings for Starlink describe a lower-altitude shell at \(\sim\!550\) km \cite{fcc2018starlinkmod}. A second analytic lineage derives from Streets-of-Coverage (SoC) polar bands and their descendants, which trade uniform \gls{raan} spacing for optimal band phasing and multiplicity structure \cite{lang1998comp}. Eccentric-orbit constellations extend these ideas beyond circular shells, illustrating the geometric richness of noncircular constellations \cite{lang1998comp}.
    
    Beyond these families, Flower Constellations (FCs) generalize the design space by enforcing compatibility with a chosen rotating frame so that all spacecraft traverse the same closed relative path; integer parameters regulate resonance and phasing, creating elegant repeating ground tracks and ``secondary paths'' that expand slotting options without sacrificing symmetry \cite{mortari2004,wilkins2008partI,wilkins2008partII,mortari2006acta}. FCs have been proposed for regional and global coverage, interferometry, and relay applications, and they unify several classical patterns under a common phasing framework \cite{mortari2004,mortari2006acta}.
    
    Optimization over these families - and over arbitrary constellations - has followed the rise of metaheuristics and multiobjective design. Genetic algorithms (GAs) have been used to search mixed discrete-continuous spaces for coverage and revisit performance \cite{ely1999GA,williams2001MOGA}. Subsequent studies adopted elitist, Pareto-based machinery (e.g., \gls{nsgaii}-style selection) to explore tradeoffs among revisit metrics and constellation parameters in a single run, producing well-spread fronts rather than a single ``best'' constellation \cite{williams2001MOGA,ferringer2006jsr,deb2002nsgaii}. Parallel processing and diversity maintenance can reduce wall-clock time for expensive constellation evaluations while preserving multiple design alternatives \cite{ferringer2007jsr}. In parallel, semi-analytical revisit/coverage estimators can substantially reduce evaluation cost and are well-suited for early design trade studies and optimization loops \cite{crisp2018revisit}. Mixed-integer formulations have also been developed for complex regional/temporal requirements, providing exact or bounded-optimal solutions when problem sizes are moderate \cite{lee2019ilp}.
    
    Applications and objective functions have broadened as mega-constellations matured. Communications studies analyze coverage probability and achievable rates at network scale using stochastic-geometry surrogates \cite{okati2020stochgeom}, and complement these models with architecture-level capacity/cost trade studies for megaconstellations \cite{delportillo2023comms}. For Earth-observation missions, common performance measures include revisit time and response time \cite{SahSrivastavaDas2021}, and constellation architectures may also be designed with on-orbit reconfiguration and slotting constraints in mind \cite{ArnasLinares2021}. Many constellation-architecture studies evaluate nominal coverage/revisit proxies under deterministic orbit assumptions \cite{ely1999GA,williams2001MOGA,ferringer2006jsr}; in operations, however, planning and safety can be sensitive to exogenous uncertainty in orbit prediction (e.g., atmospheric-density-driven drag uncertainty in LEO) and to stochastic task feasibility (e.g., cloud cover), motivating scenario-based and uncertainty-aware formulations \cite{EddyKochenderfer2019,GondelachLinares2019,NormanRivest2024}. This motivates a closer look at uncertainty-aware formulations.
    
    The synthesis for this dissertation follows a pragmatic pattern. Analytic shells provide baselines and seeds; population metaheuristics supply global exploration; and uncertainty is elevated to a first-class design driver. Beta/\gls{hpd} statistics control adaptive search, while exact finite-bank event fronts produce \(q=0.9\) maneuver--mass coverage boxes for canonical reporting. Debris surrogates remain optional stress-test tools.

    \section{Stochastic Optimization} \label{se:stochasticoptim}
    
    \textbf{Formulations for uncertainty in the problem (Definition~2).}
    When randomness enters the objective/constraints, three families are common. 
    (i) Risk-neutral stochastic programs minimize expected cost, e.g., two-stage models with recourse. 
    (ii) Chance-constrained programs bound the probability of violation, $\mathbb{P}\{g(x,\xi)\le 0\}\!\ge\!1-\varepsilon$ \cite{LaguelMalickVanAckooij2021,charnes1959chance}. 
    (iii) Risk-measure models replace expectations by coherent risk measures; \gls{cvar} is especially attractive due to convexity and sample-based estimation \cite{rockafellar2000cvar}. 
    In this work, these constructs serve as background and optional diagnostics. The baseline evaluates catalogue-derived unsafe events with Beta/\gls{hpd} adaptive control, then reconstructs exact event-count coverage fronts from maneuver and released-mass alternatives (\cref{se:optsetup}).
    
    \textbf{Sample-average approximation and scenario methods.}
    Monte Carlo \gls{saa} replaces expectations with empirical averages over i.i.d.\ scenarios; under mild conditions, \gls{saa} solutions converge with standard finite-sample validation tools \cite{AhmedShapiro2002SAA}. In this dissertation, \gls{saa} informs the use of fixed scenario banks and common random numbers (CRN) for fair benchmarking, but the baseline aggregation remains tail-oriented rather than sample-mean.
    
    \textbf{Robust and distributionally robust optimization.}
    Set-based robust optimization (RO) enforces constraints for all realizations in an uncertainty set, with tractable frameworks that allow explicit tuning of conservatism \cite{BertsimasSim2004,BeyerSendhoff2007}. Beyond fixed sets, distributionally robust optimization (DRO) protects against all distributions within an ambiguity set constructed from data or prior information. Wasserstein-ball DRO \cite{EsfahaniKuhn2018} and general convex DRO frameworks \cite{WiesemannKuhnSim2014} provide principled middle grounds between pure \gls{saa} and fully robust modeling. In this dissertation, DRO-style ablations help gauge sensitivity to surrogate misspecification.
    
    \textbf{Noisy multiobjective search and uncertainty-aware selection (Definition~1 meets Definition~2).}
    Randomized/metaheuristic solvers remain a workhorse for large, nonsmooth mission-design spaces. Under noise and scenario variability, selection noise can dominate progress unless explicitly managed. Established remedies include resampling/averaging to reduce evaluation noise and robustness-aware selection based on summary statistics of repeated evaluations \cite{BeyerSendhoff2007,deb2006robustMO}. The implementation mirrors this literature: \gls{nsgaii}-style survival selection with real-coded operators, DE mutation for continuous geometry, and \gls{pso}/SA for rapid prospecting/polishing - each wrapped around fixed-scenario evaluations with \gls{crn} and risk-aware aggregation to stabilize Pareto sorting \cite{deb2002nsgaii,StornPrice1997,coello2002mopso,Spall2003}. The algorithmic stochasticity (Definition~1) complements the problem's exogenous randomness (Definition~2).
    
    \textbf{Implications for constellation design under debris uncertainty.}
    Deterministic coverage/revisit optimization is common in constellation-architecture studies \cite{ely1999GA,williams2001MOGA,ferringer2006jsr}, while safety-critical orbit operations explicitly account for uncertainty through probabilistic and robust formulations. The adopted methodology combines fixed banks and \gls{crn}, Beta/\gls{hpd} adaptive diagnostics, exact finite-bank coverage fronts, and sensitivity framing while remaining compatible with mixed discrete-continuous metaheuristics.

    \section{Orbital Debris \& JCA} \label{se:debrisjca}
    Orbital debris poses a fundamentally statistical challenge when assessing conjunction risk. The task of cataloging and precisely tracking all hazardous objects becomes increasingly difficult as object size decreases \citep{ESA,JARRY2019185,Phipps2011lodr}, but agencies like NASA and ESA have characterized the debris environment in broad strokes. ESA's Annual Space Environment Report estimates (epoch Nov.\ 1, 2016) approximately 34,000 objects larger than 10 cm, about 900,000 objects in the 1-10 cm range, and \(\sim\!128\) million objects in the 1 mm-1 cm range \citep{ESA}; NASA's Orbital Debris Program Office maintains models and resources for debris environment assessment (e.g., ORDEM, LEGEND) \citep{Krisko2014ORDEM}. This population is far from uniformly distributed in altitude-inclination space, with concentrated bands associated with common operational regimes and major breakups (including sun-synchronous inclinations) \citep{ESA}. The proliferation of small debris is especially problematic because such fragments are typically untrackable yet can still be damaging \citep{JARRY2019185,Phipps2011lodr}. In short, the debris-on-debris collision hazard (so-called ``endogenous'' collisions) grows over time as the debris population grows, following the cascade dynamics outlined by \citet{Kessler1978}. These observations motivate the modeling choices and uncertainty-aware optimization framework developed in \cref{ch:methods}.

    Just-in-time collision avoidance (JCA) refers to rapid-response methods that impart a small $\Delta V$ (or equivalently a small perturbation) to one of the threatening objects when a conjunction is predicted, thereby preventing the collision without necessarily removing the object from orbit \citep{JARRY2019185}. The key idea is to act ``just in time'' - i.e., after a high-risk conjunction is identified but sufficiently before the projected impact - so that a modest perturbation produces a meaningful miss distance \citep{JARRY2019185,GonzaloColombo2020}. JCA is especially relevant for conjunctions involving non-maneuverable debris objects (or derelict spacecraft) that cannot execute avoidance maneuvers \citep{JARRY2019185}. The 2009 Iridium 33-Cosmos 2251 collision is a case in point: the unintentional hypervelocity impact generated debris clouds whose properties and long-term consequences have been analyzed in detail \citep{AnzMeadorLiou2010,MasonStuplMarshallLevit2011,ESA}.
    
    A number of collision-mitigation architectures have been proposed for debris-debris risk reduction. One class uses ballistically deployed particle clouds to induce a brief drag impulse \citep{JARRY2019185}. Another class applies contactless engagements such as ground-based lasers that use photon pressure to perturb debris orbits and reduce the risk of upcoming conjunctions \citep{MasonStuplMarshallLevit2011}, and related laser-based debris-removal concepts have been studied for mitigating collisional cascading in LEO \citep{Phipps2011lodr}. Across architectures, key trades include warning time, achievable impulse, and operational/legal constraints.
    
    This dissertation focuses on a dust-based JCA architecture and treats constellation sizing, placement, and operating policy as optimization variables. The operational baseline is catalogue-first, uses Beta/\gls{hpd} search diagnostics, and publishes exact finite-bank coverage-box resource fronts; statistical surrogates remain auxiliary stress-test tools.

    \medskip
    \noindent\textbf{Summary and gaps.}
    The foregoing review motivates three foci that this dissertation addresses:
    (1)~constellation architecture studies often evaluate deterministic coverage/revisit objectives, while operational performance can be sensitive to orbit-prediction and environmental uncertainty \citep{ely1999GA,williams2001MOGA,ferringer2006jsr,GondelachLinares2019,EddyKochenderfer2019,NormanRivest2024};
    (2)~stochastic and robust optimization tools - \gls{saa}, chance constraints, \gls{cvar}, and \gls{dro} - provide a principled way to encode reliability and collision-risk constraints, and are increasingly applied in conjunction assessment and collision avoidance \citep{LaguelMalickVanAckooij2021,GonzaloColombo2020,RyuBouvierLalani2024};
    (3)~particle-cloud JCA concepts have been analyzed at the maneuver/action level \citep{JARRY2019185}, and this dissertation extends the discussion to system-level constellation sizing, placement, and operating policy optimization.
    The methodology developed in \cref{ch:methods} directly targets these gaps through catalogue-derived event evidence, beta-distribution reliability semantics, and uncertainty-aware objectives/constraints (\cref{se:objectives,se:constraints}), evaluated by robust multiobjective optimizers (\cref{se:optimalgos}); debris surrogates remain optional for ablations.
